\justifying
\NSFCBodyText

\subsubsection{研究背景与意义}

跨境物流监测正在由“位置可见”转向“货物状态与风险可解释”。物流物联网、智能运输和港口数字化研究表明，传感器、边缘计算与智能分析能够提高过程透明度，但异构数据、动态链路和跨主体协同仍是落地中的持续约束\cite{golpira2021liot,rejeb2020iotlogistics,chung2021smartlogistics,digitaltwinsfor2023}。本项目面向可能经历新疆口岸至中亚长距离运输、覆盖间歇和带宽波动的应用条件，研究车辆端带宽、算力、能量及缓存共同受限时的异常证据保全。冲击、倾斜、封志开启和包装破损需要运输阶段多模态相互印证，温湿度、烟雾、门禁与视频则要在仓储阶段形成连续风险判断。冷链监测与物联网风险研究也显示，单点传感或事后追溯难以同时满足状态感知和风险管理要求\cite{badiamelis2018coldchain,tsang2018iotrisk}。

本项目关注的基本矛盾是：异常决策要求连续证据，而弱网使事件发生、设备采样与数据到达彼此错位；资源受限又迫使系统在证据尚不完整时决定“采什么、保什么、先传什么”。这种错位具有状态相关性，不能按普通随机丢失处理。当剧烈振动触发高频采样、低电量压低采样频率、弱链路延迟图像送达时，“可见数据”本身已受到事件和系统状态选择。若仍把缺失位置当作无信息的缺失掩码，模型可能把“没有看到异常”误解为“异常没有发生”；若只追求平均压缩率或吞吐量，又可能优先发送大量冗余正常片段而丢失少量关键原始证据。提高单一识别模型精度、增加缓存或提高压缩率，均不足以解决该问题。需要回答的是：决策时刻已经到达的证据在何种条件下足以辨识事件风险，资源动作如何安全改变这一证据状态，以及仓储本地风险在环境漂移后如何保持校准。厘清上述机制，有助于建立弱网环境下“可辨识—可保全—可校准”的异常证据理论，也为跨境物流全过程监测、可靠上传和分级预警提供方法基础。

\subsubsection{国内外研究现状及发展动态分析}

\textbf{（1）异步、缺失多模态异常感知。} 不规则时序建模已从插值和缺失掩码扩展到连续时间状态模型，GRU-D、潜在常微分方程等工作为不规则采样与缺失数据学习提供了基础\cite{che2018grud,rubanova2019latentode}；异步多变量异常检测进一步考虑变量间时间错位和异常定位\cite{practicalapproachto2021,aaadasynchronousinter2025,tuli2022tranad}。在视觉与传感融合方面，缺失模态恢复、工业多模态异常检测和真实场景基准持续发展\cite{distributionconsistentmodal2023,mfganmultimodalfusion2024,m3dmnrrgb2024,robustada2025,realiadd32025}。这些工作主要把缺失作为输入现象，较少区分“设备未采样”和“已经采样但尚未送达”；当缺失概率受网络、能量或事件状态影响时，普通缺失掩码与时间编码难以给出风险何时仍可辨识、何时应拒判的边界。

\textbf{（2）任务导向通信与边缘资源协同。} 语义/任务导向通信从传输比特转向保留下游任务信息，DeepSC 及其后续任务通信研究验证了任务损失驱动的传输思路\cite{xie2021deepsc,taskorientedsemantic2023,distributedtaskoriented2024,goalorientedsemantic2024}。延迟容忍网络、适应性采样、分裂推理和资源分配分别处理间歇连接、能量、带宽与计算约束\cite{reviewonrecent2022,delaytolerantnetworks2025,parametricmachinelearning2024,energyawareadaptive2023,ispliteeimage2024,artificialintelligencedefined2024}。现有策略多按单包置信度、数据新鲜度、大小或即时任务损失排序，难以表达一次异常由振动片段、封志状态和图像共同证实时的互补性；压缩、缓存、续传和卸载也常被分别优化，缺少“截止期前证据链是否充分”这一统一决策量及未知异常原始证据的安全下界。

\textbf{（3）仓储多源风险与漂移校准。} 时间序列异常检测正在由静态判别转向异常—变点区分、概念漂移适应和预测不确定性估计\cite{anomalyandchange2023,conceptdriftadaptation2022,earlyconceptdrift2025,uncertaintyawaretime2024}。概率校准和选择性预测说明，高风险应用除分类正确外，还需在证据不足时暴露不确定性或拒绝判断\cite{guo2017calibration,geifman2017selective}。然而，仓储温湿度缓变与烟雾、门禁、视频突发事件具有不同时间尺度，静态阈值难以在仓库、季节和设备变化下同时维持判别性能、校准程度与预警提前量。对完成货物级配对的冲击/倾斜或非法开启事件，入库后是否还存在事件特异的附加预测信息，也需要与仓储本地风险分开检验。

\subsubsection{现有研究不足}

现有研究仍有三处相互关联的缺口。其一，事件过程与观测/到达过程未被充分解耦，无法在状态相关缺失下区分“未发生、未采到、未送达”，完整轨迹重建也容易掩盖不可辨识性。其二，资源调度缺少相对于当前证据集合、同时考虑互补、冗余和迟到失效的条件价值；平均精度或吞吐改善不能保证关键证据得到保全。其三，仓储多源风险缺少漂移条件下的校准、拒判和可检验响应规则；运输事件摘要能否增加预测信息，则尚需在实体配对条件具备时单独验证。三处缺口共同指向一个认知问题：弱网状态相关观测下的可靠证据状态如何形成、受控改变并服务于后续风险判断。

\subsubsection{本项目的研究切入点}

本项目以“决策时刻可靠证据状态”为共同对象，沿三步展开：分离事件、采样和到达三个时钟，研究运输事件风险的条件辨识、风险界与拒判；把候选数据的价值定义为其加入现有证据集合并在截止期前交付后带来的风险降低，在关键原始证据保留、端侧最低告警、压缩保真和幂等续传约束下协调资源；建立漂移条件下可校准、可拒判的仓储本地多源风险，并在货物级配对和统计功效达到进入条件后检验运输事件摘要的附加预测价值。

据此提出三层可证伪假设：在时间元数据、可观测锚点和支持条件成立时，三时钟机制能够给出校准的运输事件风险；条件偏离时，敏感性风险区间应变宽或触发拒判。以该风险为基础的集合条件证据价值，在预算匹配且硬安全约束不变时，应降低关键异常的截止期漏警风险。仓储本地多源风险在漂移下应保持校准，并在证据不足时拒判；只有完成实体配对且达到统计进入条件，才进一步检验冲击/倾斜或非法开启摘要的条件预测增益。任一机制若只相对弱基线有效、依赖未来信息，或在条件破坏与负对照下仍给出同样结论，即被否证。

% \subsubsection{科学意义或应用前景}
% % 在此处撰写科学意义或应用前景。
